\documentclass{article}
\usepackage{amsmath, amssymb, amsthm}
\usepackage{hyperref}
\usepackage{geometry}
\geometry{margin=1in}

\title{Erd\H{o}s Problem \#311}
\date{Last updated: 2025-08-31}
\author{Source: erdosproblems.com}

\begin{document}
\maketitle

\noindent\textbf{Status:} OPEN \quad \textbf{No prize}

\noindent\textbf{Tags:} number theory, unit fractions

\medskip

\noindent\textbf{Problem Statement:}

Let $\delta(N)$ be the minimal non-zero value of $\lvert 1-\sum_{n\in A}\frac{1}{n}\rvert$ as $A$ ranges over all subsets of $\{1,\ldots,N\}$. Is it true that\[\delta(N)=e^{-(c+o(1))N}\]for some constant $c\in (0,1)$?




It is trivial that\[\delta(N)\geq \frac{1}{[1,\ldots,N]}=e^{-(1+o(1))N},\]where $[1,\ldots,N]$ is the least common multiple of $\{1,\ldots,N\}$. The formulation in \cite{ErGr80} has the additional condition that $A$ contain no $S$ such that $\sum_{n\in S}\frac{1}{n}=1$, but Kovac in the comments has shown that the simpler formulation above is equivalent. Tang has shown that\[\delta(N) \leq \exp\left(-c\frac{N}{(\log N\log\log N)^3}\right)\]for some constant $c&gt;0$.




References

[ErGr80] Erd\H{o}s, P. and Graham, R., Old and new problems and results in combinatorial number theory . Monographies de L'Enseignement Mathematique (1980).


\bigskip
\noindent\small{Source: \url{https://www.erdosproblems.com/311}}
\end{document}
